% Appendix, from exam/README.md: what the two written papers are for, what they cover, the
% conventions they assume, and how they are marked.

\chapter{The Self-Assessment Papers}
\label{exam:ch:about}

This appendix and the four after it hold two complete three-hour papers for \emph{Modern Embedded
C++}, with worked solutions. Appendices~\ref{exa:ch:paper} and \ref{exb:ch:paper} are the papers,
questions only; Appendices~\ref{ansa:ch:answers} and \ref{ansb:ch:answers} are their model answers,
with marks.

\section*{What these are for}
The book has \textbf{no exam and nothing is marked}. What you get instead is the exercise sets after
every chapter, each with a published solution, and the questions in every chapter's review.

These papers do not replace that. \textbf{They are here for you to test your own skills and
knowledge, nothing more.} They gate nothing, they are not a qualification, and no part of the book
requires them. Nothing in the companion repository depends on them, and neither \code{make build}
nor \code{make format-check} knows they exist.

What they check is something the exercises cannot. Every exercise is done at a keyboard, with the
chapter open beside you and a compiler ready to say no. These papers ask what you can reconstruct
with none of that in front of you --- which of the rules you have absorbed and which you have merely
been looking up.

\textbf{Take one after the last chapter.} Both papers draw on all six chapters, so sitting one
partway through examines material you have not read yet, and the result says more about how far
you have read than about what you have understood.

\section*{The two papers}
Both cover the whole book, Chapters~1 to 6, at the same weighting. They share no question. Either
can be used alone; use both as a main sitting and a resit, or in alternate years.

\begin{booktable}{@{}p{4.6em}Lp{5.4em}r@{}}
  \thead{Question} & \thead{Topic} & \thead{Chapters} & \thead{Marks} \\
  \midrule
  1 & From C to C++: linkage, keywords, vocabulary & 1 & 12 \\
  2 & Structs, drivers and object lifetime & 1, 2 & 12 \\
  3 & Classes: what belongs to the type & 2 & 12 \\
  4 & Copy, move, and what the compiler writes & 2 & 13 \\
  5 & Inheritance, interfaces and polymorphism & 3 & 13 \\
  6 & The factory, and who owns what & 4 & 13 \\
  7 & Templates and compile-time programming & 5 & 12 \\
  8 & Threads and synchronization & 6 & 13 \\
  \midrule
    & & & \textbf{100} \\
\end{booktable}

\textbf{Both papers mix theory with code}, because the subject does. Roughly half the marks are for
prose --- what a keyword guarantees, why an interface is worth its vtable, what a data race is ---
and roughly half are for reading or writing C++ on paper.

The two papers differ in which half they lean on:
\begin{itemize}
  \item \textbf{Paper A leans towards code that is put in front of you.} Seven of its questions hand
    you a snippet that compiles, passes a smoke test, and is wrong: a software timer with four
    defects against the book's rules, a driver missing \code{explicit}, a class whose copy
    constructor shares a heap block, an interface without a virtual destructor, a factory whose
    product is leaked, a bit-clearing template that wipes the upper half of a 64-bit register, and a
    producer/consumer that holds a mutex and still races. Naming the defect is half the marks;
    saying what it does to the running program is the other half.
  \item \textbf{Paper B leans towards code you write and design decisions you justify.} It asks for a
    C driver rewritten as a C++ class, an enumeration class and its validity check, a move
    constructor, a serial interface with every convention \secref{c3:app:b} imposes on one, a stub
    factory returning \code{std::unique_ptr}, and a fixed-capacity container with its size in the
    type.
\end{itemize}

Neither paper asks anything of a compiler that a careful reader could not answer from the chapters.
Where a question turns on a rule that is easy to state and easy to get backwards --- which
\code{const} goes in the header and which in the source file, what a user-declared destructor does
to the implicit move operations, why \code{std::move} moves nothing --- the solution says so at
length, because that is where the marks and the understanding both are.

\section*{Conventions the papers assume}
Both papers state these in their own rubric, so a candidate never has to have read this appendix.
\begin{itemize}
  \item \textbf{C++17 throughout.} Standard headers may be assumed included unless a question asks
    for them.
  \item \textbf{Code is marked on semantics, not syntax.} A missing semicolon, a forgotten
    \code{#include} or a brace in the wrong column costs nothing. A missing \code{virtual} on a
    destructor costs everything.
  \item \textbf{The book's house style is expected} where a question asks for it: \code{camelCase},
    type names with a leading capital, \code{my} on private members and \code{our} on private
    statics, a leading capital on a \code{static constexpr} member, brace initialization,
    \code{noexcept} on driver methods, \code{[[nodiscard]]} on queries and repeated on overriding
    methods, and copy and move deleted on anything that stands for a unique hardware resource.
  \item \textbf{``State the consequence'' means the consequence.} Naming a defect earns half the
    marks; saying what it does to the running program earns the other half. ``This is wrong''
    scores nothing.
  \item \textbf{Where a question asks for a stated number of defects}, listing more is not
    penalised, but only that many are marked, so the strongest answers go first.
\end{itemize}

\section*{Marking}
Every solution is written to be marked by somebody who has read the chapters and does not otherwise
write C++ for a living, so each carries the reasoning rather than the answer alone. Marks are shown
per part.

Three conventions worth agreeing before a paper is marked:
\begin{itemize}
  \item \textbf{Method carries the marks.} A correct diagnosis with a clumsy fix is worth more than a
    correct fix with no diagnosis. Several questions feed the next part, so an error should be
    followed through rather than penalised twice.
  \item \textbf{The named traps are worth full marks on their own.} A number of parts exist entirely
    to see whether a candidate walks into a specific mistake: a shallow copy constructor, a move
    assignment that forgets to release what it already owns, a non-virtual destructor in an
    interface, a discarded owning pointer, \code{1U << bit} on a 64-bit register, a mutex that one
    of the two threads never locks, a \code{std::thread} that is neither joined nor detached. Where
    a solution flags one of these, a candidate who walks into it loses those marks and no others.
  \item \textbf{The discussion parts are not decoration.} ``State why'' is where the book's actual
    content is, and a paper marked only on whether the code compiles would pass a candidate who has
    understood none of it.
\end{itemize}
